\hypertarget{phase-to-intensity-conversion-of-the-decorrelating-phase-noise}{%
\chapter{Phase-to-intensity conversion of the decorrelating phase
noise}\label{pm_am-phase-to-intensity-conversion-of-the-decorrelating-phase-noise}}

Follow-up to \texttt{docs/row\_coherence.md}, which replaces per-chip
frequency shifters by common Gaussian phase noise (3 rad rms,
correlation rate B = 2 GHz) on the seed. Code:
\texttt{scripts/m16\_pm\_am.py}; results
\texttt{results/m16\_pm\_am.json} (+ .png). Values are SIM, valid for
the element models and parameter ranges stated (CHOICE ranges, not
measured components), and should be reviewed before they are quoted.

\hypertarget{1-result}{%
\section{Result}\label{pm_am-1-result}}

Phase noise alone leaves the intensity constant: the simulated control
gives -376 dB/Hz, the numerical floor. Intensity noise appears only
where an element filters the spectrum. Two elements matter: narrow
band-pass filters and multipath (any delayed copy, from two
reflections). Gain tilt, fibre dispersion and polarisation crosstalk in
PM fibre are harmless. The criterion is a parametric heating rate Gamma
= pi\^{}2 nu\^{}2 S(2 nu) of at most 0.1 s\^{}-1 at nu = 84 kHz, i.e. a
relative intensity noise (RIN) density at 168 kHz of S \textless=
1.4e-12/Hz (-118 dB/Hz). A representative chain gives -134.5 dB/Hz
(Gamma = 2.5e-3 s\^{}-1), 16 dB inside the criterion. The design rules
that follow are: (i) no filter narrower than about 100 pm (20 times the
1.35 GHz rms modulation width) in the light path; (ii) the product of
any two reflections in the chain, R1 R2, below 1.7e-3 (amplitude r
\textless= 0.042), i.e. no pair of uncoated glass surfaces in series.

\hypertarget{2-elements}{%
\section{Elements}\label{pm_am-2-elements}}

RIN density in the 100-300 kHz band, record 131 us at 16 GS/s, worse of
two relative phases for copies:

{\footnotesize
\begin{longtable}[]{@{}>{\raggedright\arraybackslash}p{0.293\dimexpr(\linewidth-10\tabcolsep)\relax}>{\raggedright\arraybackslash}p{0.223\dimexpr(\linewidth-10\tabcolsep)\relax}>{\raggedright\arraybackslash}p{0.210\dimexpr(\linewidth-10\tabcolsep)\relax}>{\raggedright\arraybackslash}p{0.153\dimexpr(\linewidth-10\tabcolsep)\relax}>{\raggedright\arraybackslash}p{0.121\dimexpr(\linewidth-10\tabcolsep)\relax}@{}}
\toprule\noalign{}
element & case & RIN rms (all frequencies) & RIN at 168 kHz & heating
rate \\
\midrule\noalign{}
\endhead
\bottomrule\noalign{}
\endlastfoot
amplifier gain tilt & 0.05 / 0.5 / 5 dB/nm & 6e-5 / 6e-4 / 6e-3 & -237 /
-216 / -182 dB/Hz & \textless{} 1e-6 s\^{}-1 \\
band-pass, FWHM 1 nm, offset 250 pm & ASE filter & 7e-3 & -173 dB/Hz &
\textless{} 1e-6 s\^{}-1 \\
band-pass, FWHM 100 pm, centred / offset 25 pm & & 1.0e-2 / 7.0e-2 &
-130 / -133 dB/Hz & 7e-3 / 3e-3 s\^{}-1 \\
band-pass, FWHM 30 pm, centred / offset 7.5 pm & & 9.6e-2 / 0.21 &
\textbf{-110 / -113 dB/Hz} & \textbf{0.70 / 0.37 s\^{}-1} \\
multipath, R1 = R2 = 1e-2, delay \textgreater= 0.3 ns & two -20 dB
reflections & 1.4e-2 & -130 dB/Hz & 7e-3 s\^{}-1 \\
multipath, R1 = R2 = 1e-2, delay 0.1 / 0.01 ns & on chip, mm scale &
7e-3 / 1e-4 & -133 / -171 dB/Hz & 3e-3 / 1e-6 s\^{}-1 \\
multipath, R1 = R2 = 1e-3 or 1e-4, any delay & & \textless= 1.4e-3 &
\textless= -149 dB/Hz & \textless= 8e-5 s\^{}-1 \\
PM fibre, PER -20 dB, 5 / 20 m (DGD 1.2 ps/m) & then the chip TE guide
as polariser & 3e-4 / 5e-3 & -161 / -137 dB/Hz & 6e-6 / 1.4e-3
s\^{}-1 \\
dispersion, 20 ps\^{}2/km, 5 m to 1 km & & \textless= 8e-4 & \textless{}
-280 dB/Hz & 0 \\
chain (Sec. 3) & & 6.9e-3 & -134.5 dB/Hz & 2.5e-3 s\^{}-1 \\
\end{longtable}
}

The structure follows from the physics. Gain tilt and dispersion respond
to the instantaneous frequency, the time derivative of the phase, whose
spectrum vanishes towards low frequency. So their intensity noise lies
at GHz, not at 168 kHz. A delayed copy of amplitude r produces an
intensity term 2 r cos(phi(t) - phi(t - tau)). Once tau exceeds the
correlation time (about 0.3 ns), this term is a random telegraph-like
signal of variance 2 r\^{}2 spread over about B, so S \textasciitilde{}
4 r\^{}2 x 0.40 / B (the 0.40/B is the low-frequency density found in
module 15). For tau far below the correlation time the term is quadratic
in the phase increment and much smaller. A band-pass narrower than the
spectrum is the multipath effect again: it is a delay distributed over
its impulse response, about 1/(pi x FWHM) = 40 ps for 30 pm.

\textbf{Dependence on B} (r = 1e-2): at tau = 1 ns, S is -125, -130 and
-138 dB/Hz for B = 0.5, 2 and 8 GHz, because the noise spreads as 1/B.
At tau = 0.1 ns, B = 0.5 GHz is the quietest (-170 dB/Hz): there the
copy is still inside the correlation time. A higher B therefore helps
against long-delay reflections and hurts against short ones. B = 2 GHz
is a reasonable middle.

\hypertarget{3-representative-chain}{%
\section{Representative chain}\label{pm_am-3-representative-chain}}

Amplifier tilt 0.5 dB/nm; ASE filter of 1 nm FWHM, offset 250 pm; PM
fibre 5 m with -20 dB polarisation crosstalk; two -30 dB connector
reflections 5 m apart (tau = 25 ns); an on-chip etalon between a -30 dB
input facet and the 2.3\% glued emitter facet (module 14), 190 ps round
trip; 5 m of dispersion. Result: RIN rms 6.9e-3 (at GHz), -134.5 dB/Hz
at 168 kHz, Gamma = 2.5e-3 s\^{}-1 (energy e-folding about 7 min). The
on-chip etalon and the connector pair dominate.

\hypertarget{4-comparison-with-no-modulation}{%
\section{Comparison with no
modulation}\label{pm_am-4-comparison-with-no-modulation}}

Without phase modulation, the same reflections do not produce noise at
168 kHz. Instead they give a slow intensity offset 2 r cos(Phi) that
follows the drift of the path phase Phi: 2\% for r = 1e-2, at Hz to kHz
rates. The modulation converts this slow, uncontrolled offset into
broadband noise of the same rms. Spread over GHz, it deposits only 1e-13
to 1e-12 /Hz near the trap frequencies. Phase noise therefore helps
against parasitic etalons as well as against row crosstalk, provided
rule (ii) holds.

\hypertarget{5-limits-and-what-to-measure}{%
\section{Limits and what to
measure}\label{pm_am-5-limits-and-what-to-measure}}

The elements are linear and treated separately. Gain dynamics of the
amplifier (which can suppress or shape slow intensity noise) and
nonlinear effects (self-phase modulation, stimulated Brillouin
scattering) are not modelled. The criterion uses the radial frequency;
the axial frequency is lower and its parametric resonance lies where the
RIN density is the same, so the conclusion holds. Measurement: delayed
self-heterodyne or direct RIN of the modulated, amplified light after
the full chain, with the modulation on and off, at 10 kHz to 1 MHz. The
prediction to test is a floor below -130 dB/Hz with a 1 nm ASE filter
and AR-coated components.

\hypertarget{verification-notes}{%
\section{Verification notes}\label{pm_am-verification-notes}}

{\footnotesize
\begin{longtable}[]{@{}>{\raggedright\arraybackslash}p{0.265\dimexpr(\linewidth-8\tabcolsep)\relax}>{\raggedright\arraybackslash}p{0.146\dimexpr(\linewidth-8\tabcolsep)\relax}>{\raggedright\arraybackslash}p{0.146\dimexpr(\linewidth-8\tabcolsep)\relax}>{\raggedright\arraybackslash}p{0.444\dimexpr(\linewidth-8\tabcolsep)\relax}@{}}
\toprule\noalign{}
value & location & type & pointer \\
\midrule\noalign{}
\endhead
\bottomrule\noalign{}
\endlastfoot
all RIN densities, heating rates & Sec. 1-4 & SIM &
results/m16\_pm\_am.json (review) \\
element parameter ranges & Sec. 2, 3 & CHOICE & stated ranges, not
component data \\
PM-fibre DGD 1.2 ps/m, beta2 20 ps\^{}2/km & Sec. 2 & RECALLED & verify
for the chosen fibre \\
Gamma = pi\^{}2 nu\^{}2 S(2 nu) & Sec. 1 & SOURCE &
O\textquotesingle Hara et al., arXiv:physics/0003049 (module 3) \\
0.1 s\^{}-1 criterion & Sec. 1 & CHOICE & energy doubling time
\textasciitilde7 s \\
\end{longtable}
}
